In a typical Persian architecture, on top of south side windows there is a
structure called ``Tabeshband'' (shader), which controls sunlight in summer
and winter. In summer when the Sun is high, Tabeshband prevents sunlight to
enter rooms and keeps inside cooler. In the modern architecture it is verified
that the Tabeshband saves about 20\% of energy cost. Figure (1) shows a
vertical section of this design at latitude of $36^\circ\!.0$\,N with window
and Tabeshband.

Using the parameters given in the figure, calculate the maximum width of the
Tabeshband, ``$x$'', and maximum height of the window, ``$h$'', in such a way
that:
\begin{description}
\item[(i)] No direct sunlight can enter to the room in the summer solstice at
  noon.
\item[(ii)] The direct sunlight reaches the end of the room (indicated by the
  point A in the figure) in the winter solstice at noon.
\end{description}

\ioaafig{2009_TH_P4-f1.pdf}
